% latexmk -pdf presentation.tex
% LTeX: language=en-US

\documentclass[aspectratio=169]{beamer}
\usepackage{diagbox}
\usepackage{hyperref}
\usetheme[background=dark, progressbar=frametitle, titleformat=smallcaps]{metropolis}
\usepackage{csquotes}
\usepackage{pdbf}
\usepackage{biblatex}
\addbibresource{bibliography.bib}


\dbSQLText{CREATE TABLE HelloWorld(id int, text varchar); INSERT INTO HelloWorld VALUES(10, 'Hello World!'); INSERT INTO HelloWorld VALUES(120, 'Hallo Welt!');  INSERT INTO HelloWorld VALUES(100, 'bonjour le monde!');  INSERT INTO HelloWorld VALUES(42, 'Hello World!');}
\dbSQLText{CREATE TABLE test(a int, b int, [attribute with space] int); INSERT INTO test VALUES(1, 1, 1); INSERT INTO test VALUES(2, 2, 4); INSERT INTO test VALUES(3, 3, 8); INSERT INTO test VALUES(4, 4, 16);}


\title{Reproducing a Paper about a reproducibility Tool}
\subtitle{Scientific Practice}
\date{16.06.2025}

\setbeamertemplate{footline}{
  \leavevmode%
  \begin{beamercolorbox}[wd=\paperwidth,ht=2.5ex,dp=1.125ex]{author in head/foot}%
    \hspace{1em}\insertshorttitle{} -- \insertshortsubtitle{}%
    \hfill%
    \makebox[3em][r]{\insertframenumber\hspace{1em}}%
  \end{beamercolorbox}%
}

\AtBeginSection[]{}

\begin{document}

{
\setbeamertemplate{footline}{}
\frame{\titlepage}
}

\begin{frame}{Reproduction}
    If I have the same code, the same environment, the same datasets and hardware I can reproduce the results described in a paper\footnote{Please do not ask me about Machine Learning. GPUs sometimes do non-deterministic stuff.}.
\end{frame}

% \begin{frame}{Archival Work}
%     \begin{itemize}
%         \item We talked about acid free paper
%         \item But what is the \textquote{acid free paper of the digital world?}
%         \item \texttt{PDF/A} as a long archival pdf format
%         \item Virtual Machines as a time capsule
%     \end{itemize}
% \end{frame}

\begin{frame}{Lets Reproduce}
    \citetitle{pdbf} aims to be able to build an all-in-one PDF to reproduce research.
    \begin{itemize}
        \item Actual Paper (\LaTeX)
        \item Used datasets (json, sql command, csv)
        \item Interactive Diagrams
        \item Attach Virtual Machines, tarballs
        \item ...
    \end{itemize}
    Ironically they did not append their setup as a Virtual Machine Disk Image within their paper.
\end{frame}


\begin{frame}[fragile]{Reproducing the Environment to Compile PDBF}
    \small
    \begin{verbatim}
Get:1 http://deb.debian.org/debian bookworm InRelease [151 kB]                  
Get:2 http://deb.debian.org/debian bookworm-updates InRelease [55.4 kB]              
Get:3 http://deb.debian.org/debian-security bookworm-security InRelease [48.0 kB]    
Get:4 http://deb.debian.org/debian bookworm/main amd64 Packages [8793 kB]
Get:5 http://deb.debian.org/debian bookworm-updates/main amd64 Packages [512 B]
Get:6 http://deb.debian.org/debian-security bookworm-security/main amd64 Packages [265 kB]
Fetched 9313 kB in 1s (9198 kB/s)
Reading package lists...
Reading package lists...
Building dependency tree...
Reading state information...
Package phantomjs is not available, but is referred to by another package.
This may mean that the package is missing, has been obsoleted, or
is only available from another source
E: Package 'phantomjs' has no installation candidate
    \end{verbatim}
\end{frame}


\begin{frame}{Let's get \texttt{phantomjs}}
    \begin{itemize}
        \item Tried old Debian, where package server still existed $\rightarrow$ version too new
        \item Building from source failed via \texttt{aur} and direct for newer and older versions
        \item Found a precompiled binary on GitHub for a compatible version
        \item Then PDBF codebase had a bug $\rightarrow$ added a missing maven dependency $\rightarrow$ workaround $\rightarrow$ this presentation
    \end{itemize}
    
\end{frame}


\begin{frame}{Dependency problems}
    \begin{itemize}
        \item Package server may go down
        \item Used releases may be yanked for security reasons and no longer available
        \item Packages may be updated by repository maintainers e.g. apt
        \item Packages may not be available for a newer system architecture (e.g. x86\_64 or arm64v8)
        \item Virtualizing a System may become more difficult
        \item Files cannot be read because of proprietary file formats
    \end{itemize}
\end{frame}


\begin{frame}{How to archive digital work}[allowframebreaks]
    \begin{itemize}
        \item Pin all dependency versions using commit-shas or lock-files from tools like cargo, uv, npm, maven, nix ...
        \item Package not only the code and datasets BUT all of their (compiled) dependencies, compiler versions, drivers, operating systems, kernels ...
        \item E.g. compiler optimization bugs in polynomial operations caused wrong results
        \item Non-Hardware reliant projects may be much easier to package
        \item Rely on Open Source Industry Standard tools, so that Commercial interest funds these project
        \item Package an OCI Image / Docker Compose setup or Virtual Machine Disk Image
        \item If not possible have exact setup description (maybe special hardware is required e.g. FPGAs) and package as much as possible
        \item System configs may require deep setup understanding and are hard to reproduce even by the authors
        \item Tools like \citetitle{reprozip} can assist identify and package relevant components through syscall tracing
        \item Store artifact in a reliable archive like \url{https://zenodo.org/} or Code with GitHub Arctic Code Vault
        \item Test the reproducible setup yourself and hand it other people
        \item Do not rely on any external servers
        \item Try to minimize non-packable dependencies like Hypervisors or Container runtimes
    \end{itemize}
\end{frame}


\begin{frame}{Just some notable features of PDBF}
    They can package interactive graphics and databases. Just change the file ending between \texttt{html} and \texttt{pdf}

    When viewing as \texttt{html} click on the Table to play around with sql.

    \dataTable[color=white]{SELECT * FROM HelloWorld;}
\end{frame}


\begin{frame}{Just some notable features of PDBF}
    I had to apply a workaround to make it build. Sadly the pdf version has slightly broken charts.
    \sql{SELECT * FROM test;}\\[3pt]
    The table \texttt{test} contains \dataText{SELECT COUNT(*) FROM test;} tuples.\\[3pt]

    The following points are linked to the chart on the left side below: \dataText[linkTo=testname, linkSelector=2, linkLabel=Value at 2]{SELECT * FROM test WHERE a=2;}, \dataText[linkTo=testname, linkSelector=3, linkLabel=Value at 3]{SELECT * FROM test WHERE a=3;}.\\[3pt]
    \chart[width=0.45\textwidth, height=0.3\textwidth, logscale=true, chartType=line, name=testname]{SELECT a, b, [attribute with space] FROM test;}
    \chart[width=0.45\textwidth, height=0.3\textwidth, chartType=bar]{SELECT a, b, [attribute with space] FROM test;}\\[1pt]
\end{frame}


\begin{frame}{Reproducibility of this presentation}
    \begin{itemize}
    \item You will receive a PDBF file of this presentation
    \item PDBF could bundle a \texttt{tar} file, but the OCI Image is around 5~GiB and I cannot send this via Slack
    \item There is a repository, a Dockerfile and a Docker Image \url{https://github.com/fidoriel/pdbf}
    \item You should be able to use the \texttt{README.md} and \texttt{Docker} to build PDBF and this presentation
    \end{itemize}
\end{frame}


\begin{frame}{Discussion}
    Reproduction is time-consuming for all participants
    \begin{itemize}
    \item Should reproducibility be a requirement in peer reviewed magazines?
    \item Can we ask reviewers to spend extra time on reproducing results?
    \item Is it reasonable to ask scientists to spend time on reproducibility, so their work can be archived better and may have a stronger stand?
    \item How to handle research that cannot be published this way, for example Intellectual Property or Privacy rights prohibit publishing of a dataset?
    \end{itemize}
\end{frame}

\begin{frame}[allowframebreaks]

\frametitle{References}
\printbibliography

\end{frame}
\end{document}
